\begin{frame}{Ratios de solvencia}
\color{solvencia}\textbf{Estructura de financiamiento y capacidad de cumplir obligaciones}\color{black}

\vspace{0.7em}
Relacionan activos, pasivos y patrimonio. Ayudan a evaluar dependencia de terceros y riesgo financiero.
\end{frame}

\begin{frame}{Razón de endeudamiento}
\[
\text{Endeudamiento}=\frac{\text{Pasivo total}}{\text{Activo total}}\times 100
\]
\textbf{Datos necesarios:} pasivo total y activo total.
\end{frame}

\begin{frame}{Caso práctico: endeudamiento}
\casoeducativo
Industrial Callao S.A.C. registra activos por \soles{2,000,000.00} y pasivos por \soles{900,000.00}.
\[
\frac{900,000.00}{2,000,000.00}\times 100=45.00\text{ \char37}
\]
\textbf{Interpretación:} el 45.00 \char37{} de los activos se financia con obligaciones de terceros.

\textbf{Decisión:} comparar con sector, costo financiero y capacidad de pago.
\end{frame}

\begin{frame}{Relación deuda-patrimonio}
\[
\text{Deuda-patrimonio}=\frac{\text{Pasivo total}}{\text{Patrimonio}}
\]
\[
\frac{900,000.00}{1,100,000.00}=0.82
\]
\textbf{Interpretación:} por cada \soles{1.00} de patrimonio, existen \soles{0.82} financiados por terceros.
\end{frame}
